\begin{textsample}{NEG Dim 2 – persona\_gpt – Score 5.00 – t982\_gpt.txt}  \label{ex:f2_neg_050}
Uniqlo has joined the growing list of global \textbf{fashion} \textbf{brands} committing to eliminate all releases of hazardous chemicals from their global supply chain and products by 2020.
This pledge from Fast Retailing Group, Uniqlo ’s parent company, follows similar commitments already made by Zara, Mango, Esprit, and Levi ’s.
As part of this commitment, Fast Retailing will disclose discharge data from 80 % of its suppliers in 2013.
The company has also committed to help lead the development of alternatives to hazardous chemicals, pushing the industry toward safer production methods.
This promise extends beyond Uniqlo to other Fast Retailing \textbf{brands}, including Comptoir des Cotonniers and Theory.
These steps are intended to benefit communities living near factories and waterways affected by toxic pollution from clothing production, as well as consumers who buy and wear these products.
By making discharge data public and working on safer alternatives, Uniqlo and its sister \textbf{brands} are taking concrete steps toward greater transparency and a toxic-free future for \textbf{fashion}.
However, many major \textbf{brands} are still lagging behind.
Companies like GAP and Calvin Klein have yet to match this level of commitment to transparent, toxic-free practices.
To help close this gap, people are being urged to share the Detox \textbf{Fashion} video and use their voices to increase pressure on these companies.
Sharing the video is presented as a way to push more \textbf{brands} to follow the lead of Uniqlo, Zara, Mango, Esprit, and Levi ’s, and to move the entire \textbf{fashion} industry toward a future free from hazardous chemicals.

% matched lemmas: brand, fashion
\end{textsample}
